\notepage{20}{Time evolution and position and energy measurements}
\smallheading{IV. Time evolution}
Classical evolution follows Hamilton's equations,
\[
 \dot x=\frac{\partial\mathcal H}{\partial p},\qquad
 \dot p=-\frac{\partial\mathcal H}{\partial x}.
\]
Quantum evolution follows the Schr\"odinger equation,
\[
 \boxed{\ii\hbar\frac{\dd}{\dd t}\ket{\psi(t)}=H\ket{\psi(t)}.}
\]
Here $H(X,P)$ is the Hamiltonian operator corresponding to
$\mathcal H(x,p)$, with an appropriate ordering of noncommuting
factors when necessary.

\smallheading{States and normalization}
Use $\braket{\psi}{\psi}=1$ for a given physical state. If a state is
constructed as $\ket{\psi}=\Omega\ket{\phi}$, its norm need not be 1;
when it is nonzero, replace it by
$\ket{\psi}/\sqrt{\braket{\psi}{\psi}}$.
The position basis is complete and delta normalized:
\[
 \braket{x'}{x}=\delta(x-x'),\qquad
 \int\ket{x}\bra{x}\dd x=I.
\]
In three spatial dimensions, write
$\ket{x,y,z}=\ket{\mathbf r}$. Then
\[
 \braket{\mathbf r'}{\mathbf r}
 =\delta^{(3)}(\mathbf r-\mathbf r'),\qquad
 \int\ket{\mathbf r}\bra{\mathbf r}\dd^3r=I,
\]
\[
 \ket{\psi}=\int\ket{\mathbf r}\braket{\mathbf r}{\psi}\dd^3r
 =\int\psi(\mathbf r)\ket{\mathbf r}\dd^3r.
\]

\smallheading{Position measurement}
Measuring $(X,Y,Z)$ gives a position in the volume element $\dd^3r$
with probability
\[
 \dd P=\abs{\braket{\mathbf r}{\psi}}^2\dd^3r
      =\abs{\psi(\mathbf r)}^2\dd^3r.
\]
An ideal exact-position outcome is represented formally by
$\ket{\mathbf r}$. Measurement is a state-selection rule, not simply
the action $\mathbf R\ket{\psi}$ of the position operator.

\smallheading{Energy measurement}
For a discrete nondegenerate spectrum,
\[
 \ket{\psi}=\sum_n\ket{E_n}\braket{E_n}{\psi},\qquad
 \phi_n(\mathbf r)=\braket{\mathbf r}{E_n}.
\]
An energy measurement gives $E_n$ with probability
\begin{align*}
 P(E_n)&=\abs{\braket{E_n}{\psi}}^2\\
 &=\abs{\int\braket{E_n}{\mathbf r}\braket{\mathbf r}{\psi}\dd^3r}^2
 =\abs{\int\phi_n(\mathbf r)^*\psi(\mathbf r)\dd^3r}^2.
\end{align*}
